%! $n$th Roots \& Rational Exponents — Unit 6, Lesson 6.1 — Exit Ticket
\documentclass[10pt]{article}
\usepackage{algebra2-article}
\usepackage{algebra2-boxes}
\newcommand{\TallMath}[1]{$\displaystyle #1\rule[-1.4em]{0pt}{3.2em}$}

\begin{document}

\pageheader{Unit 6, Lesson 6.1}{Exit Ticket}
\namedateperiod

\vspace{0.10in}

No notes. Assume all variables represent positive real numbers unless a question says otherwise.

\begin{enumerate}[leftmargin=1.8em, itemsep=10pt]
  \item \textbf{Translate and evaluate.}

        \vspace{4pt}
        (a) Write $\sqrt[5]{x^{3}}$ with a rational exponent: \blank{2.4cm}

        \vspace{4pt}
        (b) Write $y^{3/4}$ in radical form: \blank{2.8cm}

        \vspace{4pt}
        (c) $27^{2/3} = \blank{2.0cm}$ \hspace{0.6in}
        (d) $16^{-1/2} = \blank{2.0cm}$

        \vspace{4pt}
        (e) $x^{1/2}\cdot x^{1/6} = \blank{2.4cm}$

  \item \textbf{Explain.} Using rational exponents, explain why $\sqrt[3]{-64}$ \emph{is} a real
        number but $\sqrt{-64}$ is \emph{not}. Your answer must mention the \textbf{denominator} of
        the exponent.
        \writelines{3}

  \item \textbf{SOL-style multiple choice.} Which expression is equivalent to $\sqrt[4]{x^{3}}$?
        \begin{enumerate}[label=\Alph*., leftmargin=1.9em, itemsep=1pt]
          \item $x^{4/3}$
          \item $x^{3/4}$
          \item $x^{12}$
          \item $x^{-3/4}$
        \end{enumerate}
\end{enumerate}

\end{document}
